% Extension-space branch: complete proofs. AI-produced by GPT-6 Astra.
% Intended for the root paper's resultenv environment (id,status,title).
\subsection{The actual extension space, its marked points, and the cubical comparison}

Throughout this subsection, $\Sigma$ is a finite arrangement of $n$ distinct
pseudolines in $\mathbb {RP}^{2}$, any two meeting transversely at precisely one
point, and not all concurrent. Thus its oriented matroid $M$ is simple of rank
three; concurrent subfamilies are allowed. We fix orientations of the old
pseudolines only to write sign vectors. Let $\mathcal C^*(M)$ denote the set of
signed cocircuits. The extension poset $P=\mathcal E(M)$ consists of rank-three
one-element extensions by a labelled element $f$ which is not a loop. Equivalently,
these are the nonzero localization maps
\[
 \sigma:\mathcal C^*(M)\longrightarrow\{0,+,-\},\qquad
 \sigma(-Y)=-\sigma(Y).
\]
The partial order is coordinatewise, with $0<+$ and $0<-$ and with $+$ and $-$
incomparable. We write $|P|$ for the geometric realization of its order complex.
The involution $\tau(\sigma)=-\sigma$ reorients only the new element. It does not
reorient any old element and does not relabel the arrangement.

For an old element $g$, the two parallel extensions are denoted $[g^+]$ and
$[g^-]$. We call every other extension \emph{ordinary}. An ordinary extension
is represented by a pseudoline distinct from the old pseudolines. We write
$\nu(D)$ for the number of old vertices on that pseudoline, equivalently the
number of antipodal cocircuit pairs on which its localization vanishes. Set
\[
 c_j=\#\{\text{ordinary extensions modulo }\tau:\nu(D)=j\},
 \qquad F(t)=\sum_{j\geq0}c_jt^j.
\]
The old vertices form a finite set $V(\Sigma)$; $\mu(v)$ is the number of old
pseudolines through $v$, and $n_i$ is the number of vertices on the $i$th old
pseudoline. In particular $\sum_v\mu(v)=\sum_i n_i$.

\begin{resultenv}{E1}{Known (quotient consequence proved)}{The extension quotient}
The space $|P|$ is homotopy equivalent to $S^2$. The action $\tau$ is free and
there is a homotopy equivalence
\[
 |P|/\langle\tau\rangle\simeq\mathbb {RP}^{2}.
\]
More precisely, for every old element $g$ the map
\[
 \Pi_g:P\longrightarrow Q_g:=\{g^+,g^-\}\oplus\mathcal E(M/g)
\]
is a $\tau$-equivariant homotopy equivalence on order complexes. Here $\oplus$
means the ordinal sum: the two incomparable symbols $g^+,g^-$ lie below every
element of $\mathcal E(M/g)$. On ordinary non-$g$-parallel extensions,
$\Pi_g(\sigma)$ is restriction to cocircuits $Y$ with $Y_g=0$; on $[g^\pm]$
it is the corresponding symbol $g^\pm$.
\end{resultenv}
\begin{proof}
The sphere statement is Sturmfels and Ziegler's Theorem~1.2, p.~25, together
with their Corollary~4.5, p.~40; the displayed map and its homotopy equivalence
are their Proposition~4.2, pp.~37--38 \cite{SturmfelsZiegler1993}. We check the
hypotheses actually needed. The contraction $M/g$ has rank two. Every
one-element extension $M\cup f$ under consideration still has rank three,
and every rank-three oriented matroid program is Euclidean, as used in their
Corollary~4.5 and the proof of Proposition~4.2. The latter proposition therefore
applies to each choice of $g$, not merely to a special realizable case.

The formula for $\Pi_g$ commutes with $\sigma\mapsto-\sigma$: this negates the
restricted localization and exchanges the two suspension symbols. The
rank-two extension poset is the face poset of a circle, with sign reversal
acting antipodally. To see the rank-two assertion directly, represent the
rank-two oriented matroid by nonzero vectors in the oriented plane, with each
parallel class represented by one direction. The new vector's direction runs
around a circle; the finitely many old signed directions are the vertices
and the intervening intervals are the edges. Parallel elements and loops of
$M/g$ do not alter this circle. Suspending gives an actual $S^2$ with the
ordinary antipodal action: in suspension coordinates $(z,t)$ it is
$(z,t)\mapsto(-z,-t)$.

The action on $|P|$ is free. Indeed, $\sigma\neq-\sigma$ for every nonzero
localization, and an invariant finite chain under an order-preserving
involution would have every member fixed. Thus it fixes no simplex setwise
and no point. Both spaces are finite CW complexes and their quotients are
covered by these free two-sheeted covers. Since $|\Pi_g|$ is a homotopy
equivalence and both total spaces are simply connected, the induced quotient
map is an isomorphism on fundamental groups (both deck groups are $C_2$) and,
by the covering-space lifting isomorphisms, on every higher homotopy group.
The CW Whitehead theorem proves that the quotient map is a homotopy
equivalence. This establishes the quotient conclusion from the actual
specified equivariant map, rather than from an unmarked sphere homotopy type.
\end{proof}

\begin{resultenv}{E2}{Proved}{An embedded regular cubical and polygonal model}
The opposite poset $P^{\mathrm{op}}$ is the face poset of a finite regular CW
complex $T(M)$ with the following cells: an ordinary extension $D$ labels a
$\nu(D)$-cube, and each parallel extension labels a polygonal two-cell.
Every closed cell embeds. The barycentric subdivision of $T(M)$ is $|P|$,
and its quotient $\overline T(M)=T(M)/\langle\tau\rangle$ has one ordinary
cube per unoriented relative position and $n$ parallel polygonal two-cells.
Consequently
\[
 \boxed{F(-1)=1-n.}
\]
\end{resultenv}
\begin{proof}
Let $\sigma$ be ordinary, with zeros at $\nu$ old vertices. Distinct such
vertices involve disjoint sets of old pseudolines: otherwise the added
pseudoline would meet one old pseudoline at two distinct points. Choose
pairwise disjoint small disks about these vertices. Resolving the added
pseudoline to either side inside any chosen disk and leaving it fixed
elsewhere preserves its single transverse intersection with every old
pseudoline. These resolutions can be made independently. At each vertex we
may leave the passage unresolved, or choose either resolution. They give
exactly $3^\nu$ extensions above $\sigma$: no sign outside its zero set may
change in a greater localization, and its signs on the zero pairs determine
the extension. The latter uniqueness is the localization determination
theorem; it is also immediate here by gluing the specified local arcs back
to the fixed complement. With order reversed, these are precisely the
nonempty faces of $[-1,1]^\nu$, where a zero coordinate means a free cube
coordinate. Thus the proper upper interval $P_{>\sigma}$ is the face poset
of the boundary of that cube, with order reversed.

For a parallel extension $[i^\pm]$, Sturmfels--Ziegler Lemma~4.1, p.~37,
identifies $P_{>[i^\pm]}$ with $\mathcal E(M/i)$, a rank-two circle. It is
therefore a polygonal boundary, and adding its cone gives a two-cell.

We make the embedding assertion explicit. In the finite simplicial complex
$|P^{\mathrm{op}}|$, for each element $x$ take the principal-ideal
subcomplex
\[
 C_x=|\Delta((P^{\mathrm{op}})_{\leq x})|.
\]
It is literally the cone, with vertex $x$, on
$|\Delta((P^{\mathrm{op}})_{<x})|$. The preceding interval calculations say
that this is the barycentrically subdivided closed cube, or the subdivided
closed polygon, appropriate to $x$. Its open cell is $C_x$ minus that
boundary. Every point of the order complex belongs to exactly one such
open cell, as seen from the largest element of its supporting chain.
Intersections of closed cells are unions of their common closed faces.
Thus these are embedded characteristic balls in a finite regular CW
structure; no assertion that an arbitrary colimit of embeddings is embedded
is being used.

For any $x$, choose a cocircuit on which its localization is nonzero. Every
extension above $x$ has that same sign. Hence $C_x$ contains no pair of
antipodal extensions and no pair of points identified by $\tau$. The quotient
map therefore remains injective on each $C_x$ and gives the asserted regular
CW structure on $\overline T(M)$. Its Euler characteristic is
$F(-1)+n$; E1 makes this equal to $\chi(\mathbb {RP}^2)=1$.
\end{proof}

\begin{resultenv}{E3}{Proved}{The derivative identity from a fixed old vertex}
For every $v\in V(\Sigma)$,
\[
 \sum_{D\ni v}(-1)^{\nu(D)-1}=\mu(v),
\]
where the sum is over ordinary unoriented relative positions through $v$.
In particular
\[
 \boxed{F'(-1)=\sum_{v\in V(\Sigma)}\mu(v)=\sum_i n_i.}
\]
\end{resultenv}
\begin{proof}
Choose an oriented cocircuit $Y$ at $v$ and an old element $g$ through $v$.
Let $P_v=\{\sigma\in P:\sigma(Y)=0\}$, with its induced order. The condition
is independent of the choice between $Y$ and $-Y$. Since $Y_g=0$, the map
$\Pi_g$ of E1 restricts to
\[
 P_v\longrightarrow\{g^+,g^-\}\oplus R_v,
 \qquad R_v=\{\rho\in\mathcal E(M/g):\rho(Y)=0\}.
\]
In rank two, $R_v$ consists of two antipodal parallel extensions and is an
antichain with two elements. Thus the target order complex is a circle.

We verify that this restriction still induces a homotopy equivalence; this
does not follow merely by restricting an arbitrary homotopy equivalence.
Sturmfels and Ziegler prove Proposition~4.2 using their Lemma~3.2, p.~29
(Babson's fiber criterion). The proof verifies that each point fiber of
$\Pi_g$ is contractible, and that
\[
 \Pi_g^{-1}(q)\cap P_{\geq p}
\]
is contractible whenever $\Pi_g(p)<q$. For a restriction to the full inverse
image of any subposet of $Q_g$, each of these fibers and slices is unchanged.
The same criterion therefore applies to the restricted map. In the present
case one can see the slices especially explicitly: the only comparisons
in the target are $g^\pm<\rho$, and the corresponding slice is the singleton
$\{\kappa_\pm(\rho)\}$ of their Lemma~4.1. We have proved $|P_v|\simeq S^1$.
The sign action is free, so $\chi(|P_v|/\langle\tau\rangle)=0$.

In the opposite induced poset, an ordinary extension through $v$ has a
$(\nu(D)-1)$-cube as closure: the zero at $v$ must remain zero and the other
zeros can be resolved independently. A parallel extension at an old line
through $v$ has a one-cell as closure, since its proper upper interval in
$P_v$ is exactly the two rank-two localizations vanishing at the cocircuit
of $v$. There are $\mu(v)$ such one-cells modulo $\tau$. The principal-ideal
construction of E2 again proves that all these cells embed. Therefore
\[
 0=\chi(|P_v|/\langle\tau\rangle)
   =\sum_{D\ni v}(-1)^{\nu(D)-1}-\mu(v).
\]
Summing over $v$ counts an ordinary extension with $\nu(D)$ zeros exactly
$\nu(D)$ times and gives the derivative identity.
\end{proof}

\begin{resultenv}{E4}{Proved}{The marked puncture and real blow-up comparison}
Choose pairwise disjoint small open two-disks about the $n$ parallel points
in $\overline T(M)$; these points lie in genuine open two-disk neighborhoods.
Let $U(M)$ be their complement. There is a homotopy equivalence of pairs
\[
 \bigl(\mathbb {RP}^2\setminus\textstyle\bigcup_{i=1}^n\mathring B_i,
       \coprod_i\partial B_i\bigr)
 \longrightarrow
 \bigl(U(M),\coprod_i\partial B_i'\bigr)
\]
whose restrictions to the boundary circles are homeomorphisms matching the
old-line labels. Replace each removed disk by a M\"obius band, attaching its
boundary by the same circle parametrization, and denote the resulting CW
complex by $B(M)$. Then
\[
 B(M)\simeq\#^{\,n+1}\mathbb {RP}^2,
 \qquad\chi(B(M))=1-n.
\]
This statement concerns a specified operation at regular parallel points;
it does not assume that the whole extension complex is a surface.
\end{resultenv}
\begin{proof}
We give the marked argument because an unmarked homotopy equivalence to
$\mathbb {RP}^2$ alone would not justify puncturing.
Work first in the oriented complex $|P|$, and fix $g$. Put
$C=\mathcal E(M/g)$, a circle with antipodal sign action, regarded as a
face poset. Its zero-dimensional cells will be denoted $\rho$ and its
one-dimensional cells $\gamma$. For a closed edge $\gamma$ define
\[
 Q_\gamma=\{g^+,g^-\}\oplus C_{\leq\gamma},\quad
 R_\gamma=|\Pi_g^{-1}(Q_\gamma)|.
\]
For a vertex $\rho$ define
\[
 Q_\rho=\{g^+,g^-,\rho\},\qquad
 A_\rho=|\Pi_g^{-1}(Q_\rho)|.
\]
Each $|Q_\gamma|$ is the suspension of an interval, hence a closed disk;
each $|Q_\rho|$ is an interval with endpoints $g^+,g^-$. By the restricted
fiber argument in E3, every $R_\gamma$ and every $A_\rho$ is contractible.
The $R_\gamma$ cover $|P|$. Adjacent sectors meet in $A_\rho$; nonadjacent
sectors meet only in the two points $[g^+],[g^-]$. These assertions hold as
literal equalities of simplicial subcomplexes, because a chain in a circle's
face poset lies over one closed edge.

We describe the local marked neighborhoods. The star of any parallel
extension $[h]$ is a cone on the rank-two circle $\mathcal E(M/h)$, by
Lemma~4.1, and $[h]$ is minimal in $P$. Thus its star is a genuine open
two-disk neighborhood in the entire order complex. For $h\neq g$, its image
under $\Pi_g$ is the rank-two parallel vertex $\rho$ corresponding to
$v=g\cap h$. In its link $\mathcal E(M/h)$, evaluation at the cocircuit of
$v$ has two zeros and separates the circle into the positive and negative
arcs. The two zeros form the link of $[h]$ in $A_\rho$; the two closed arcs
form its links in the two adjacent $R_\gamma$. Hence $[h]$ is locally an
interior interval point of $A_\rho$ and a boundary half-disk point of each
adjacent sector. Similarly, $[g^\pm]$ is an endpoint of each $A_\rho$ and a
boundary point of each sector. These local claims use the actual rank-two
circle, not a homotopy description of a link.

We next prove that the other marked points on a meridian occur between its
two endpoints. Fix $[h]\in A_\rho$. Choose an old vertex $w\in h$, $w\neq v$.
Such a vertex exists: otherwise every old line would meet $h$ at $v$, so the
whole arrangement would be concurrent. Choose the signed cocircuit $Y_w$
with $(Y_w)_g=+$. The localization evaluation at $Y_w$, interpreted as
$-1,0,1$ and extended linearly over every order simplex, defines a continuous
map
\[
 e_w:A_\rho\longrightarrow[-1,1].
\]
A simplex cannot have both a positive and a negative value, since its
vertices are comparable localizations. The only vertex of $A_\rho$ at
which this evaluation vanishes is $[h]$: every extension in the fiber
already passes through $v$, and an ordinary pseudoline cannot also pass
through $w$ on the same old line $h$. The two parallel extensions with
underlying line $h$ have different images $\rho,-\rho$, so only one lies
in this meridian. Thus $e_w^{-1}(0)=\{[h]\}$, while its values at $[g^+]$
and $[g^-]$ have opposite signs. Each such marked point separates the two
endpoints.

Here is the elementary topological fact we use to cut a meridian. If a
contractible finite CW complex has a point with an open interval
neighborhood, cutting out a smaller interval gives exactly two contractible
components whenever the cut point separates two specified points. To prove
it, restore the interval as a one-cell joining the cut complex. Its relative
homology and the long exact sequence show that the cut complex has two
acyclic components. Van Kampen identifies the fundamental group of the
restored complex with the free product of the groups of those components,
so both are simply connected. A simply connected acyclic CW complex is
contractible, by Hurewicz followed by Whitehead. Iterating proves the
statement for finitely many disjoint cuts. Since every marked point in
$A_\rho$ separates its endpoints, their separating order is linear: cutting
one places every other point on one of the two endpoint sides. No point
can be on a branch outside the endpoint-to-endpoint order.

Choose disjoint sufficiently small polygonal disks about all $2n$ oriented
parallel points, compatible with the link sectors just described and
invariant under $\tau$. Let $P^\circ$ be the complement of their interiors,
and let $R_\gamma^\circ,A_\rho^\circ$ be the induced cut pieces. Each
$R_\gamma^\circ$ is contractible. Indeed, in $R_\gamma$ each removed piece
is a boundary half-disk (a sector at $g^\pm$ is homeomorphic to a half-disk),
and the half-disk strongly deformation retracts to its frontier semicircle
while fixing that frontier. Thus deleting it leaves a deformation retract
of the original sector. The meridian $A_\rho^\circ$ is a disjoint union of
contractible components: endpoint cuts do not change its homotopy type,
and the intervening interval cuts were handled in the previous paragraph.
The cuts at $g^\pm$ eliminate all intersections of nonadjacent sectors and
all triple intersections. All pieces and intersections can be triangulated
as finite subpolyhedra, so their inclusions are cofibrations.

We now construct the comparison map with its boundary marking. Start with
the ordinary suspension sphere $S^2=\operatorname{susp}|C|$, whose meridians
are indexed by $\rho$ and whose sectors are indexed by $\gamma$. On the
$\rho$-meridian place one marked point for each $[h]\in A_\rho$, in the
separating order just proved, and mark the two suspension tips. Remove small
disks about these $2n$ marked points. We obtain a sphere $S^\circ$ with
$2n$ boundary circles. The cut sectors are disks and the cut meridians are
intervals in exactly the same incidence diagram as
$R_\gamma^\circ,A_\rho^\circ$.

Map each boundary circle homeomorphically to the corresponding actual
circle in $P^\circ$, matching its two half-circle pieces at a non-tip
parallel point and its sector pieces at a tip. For each cut meridian
interval, its two endpoints have now been prescribed in the corresponding
contractible component of $A_\rho^\circ$; choose a path between them in
that component. These paths give compatible maps on all sector boundaries.
Each such boundary map extends over the corresponding sector disk because
$R_\gamma^\circ$ is contractible. Gluing gives a continuous map
$S^\circ\to P^\circ$ which is a homeomorphism on every marked boundary
circle. Its restriction to every sector is a homotopy equivalence between
contractible spaces, and its restriction to each intersection component is
likewise a homotopy equivalence. Gluing finite CW homotopy equivalences
along these cofibrations proves that the whole map is a homotopy
equivalence. One may verify this gluing assertion inductively by writing
each addition of a sector as a pushout along the disjoint union of its
intersection components; such a pushout is a homotopy pushout because all
inclusions are cofibrations.

All choices can be made equivariantly. The sign involution exchanges the
tips, sends $\rho$ to its antipodal vertex, and reverses each meridian's
endpoint order. On the model sphere choose the action $(z,t)\mapsto(-z,-t)$;
place paired marked points at opposite heights on antipodal meridians.
Choose the boundary maps, meridian paths, and sector fillings for one
representative of each antipodal pair and transport them by $\tau$. No
sector or meridian is invariant under this action, so these prescriptions
have no fixed-piece compatibility obstruction. The resulting equivariant
map descends to a homotopy equivalence of quotient pairs: its lift is a
homotopy equivalence, its induced map on the double-cover deck groups is
the identity, and the covering-space exact sequence and Whitehead give
the quotient conclusion. The model quotient is exactly
$\mathbb {RP}^2$ with $n$ open disks removed, not an unspecified quotient
of a homotopy sphere.

Finally attach a M\"obius band to every boundary circle on both sides of
this pair equivalence, using the same parametrizations. The boundary maps
are homeomorphisms and the attachments are cofibrations; the same homotopy
pushout argument extends the equivalence. On the model surface this is the
real point blow-up at each of $n$ points, giving the connected sum of
$n+1$ projective planes. Each replacement lowers Euler characteristic by
one, which also proves the displayed Euler formula directly.
\end{proof}

\begin{resultenv}{E5}{Proved under stated hypotheses (T3 verifies all here)}{The numerical correction forced by the cubical extension counts}
Suppose the manuscript surface has the verified cell counts
\[
 s_2=F(1),\qquad s_0=F'(1)+\sum_i n_i,
\]
and is a closed connected nonorientable surface. Then
\[
 \chi(\mathcal X)
 =1-n-4\sum_{k\geq1}k\,c_{2k+1},
 \qquad
 \mathcal X\cong\#^{\,n+1+4\sum_{k\geq1}k c_{2k+1}}\mathbb {RP}^2.
\]
It is homotopy equivalent to $B(M)$ if and only if no ordinary relative
position passes through three old vertices. This hypothesis is automatic
for $n\leq5$.
\end{resultenv}
\begin{proof}
All the stated surface hypotheses hold for $\mathcal X=\mathcal X_2$ by
T3. They are not consequences of the extension-space theorem alone.
Under them $\chi(\mathcal X)=F(1)-F'(1)-F'(-1)$ by E3. Subtract E2 and
compare coefficients. For the coefficient of $c_j$ the difference is
\[
 1-j-j(-1)^{j-1}-(-1)^j=(j-1)((-1)^j-1).
\]
This is zero for even $j$ and $-2(j-1)$ for odd $j$, proving the formula.
The classification of compact surfaces gives the homeomorphism assertion.
By E4, $B(M)$ has the homotopy type of the connected sum of $n+1$ projective
planes. These two surfaces are homotopy equivalent exactly when their Euler
characteristics agree, which here means $c_{2k+1}=0$ for every $k\geq1$.
If any ordinary position has $\nu\geq3$, resolving all but three of its
vertices produces an ordinary position with exactly three zeros, so $c_3>0$.
Conversely $c_3=0$ excludes all such positions. Finally distinct vertices on
an ordinary added pseudoline involve disjoint old-line classes, each of size
at least two; hence $2\nu\leq n$ and $\nu\geq3$ requires $n\geq6$.
\end{proof}

\begin{resultenv}{E6}{Refuted (corrected comparison proved)}{What the cyclic comparison actually selects}
Lead L2 is false: all $2^\nu$ stable generizations are present in
$\mathcal X$. Its true replacement is the following explicit cellular
description. Subdivide every ordinary cube of $\overline T(M)$ by its
coordinate midplanes. Its vertices are the centers of the original cube
faces, hence ordinary extension labels; its edges are the ordinary cover
relations. Replace each parallel polygon by the M\"obius band of E4,
subdivided as described below. In this CW model of $B(M)$ the dual
quadrangulation $\mathcal X^\vee$ is an embedded subcomplex containing
the entire one-skeleton.

Specifically, if $\sigma$ has distinct zero coordinates $i,j$ and resolving them
with specified signs gives $\sigma_i,\sigma_j,\sigma_{ij}$, its diamond is
\[
 \sigma<\sigma_i<\sigma_{ij},\qquad
 \sigma<\sigma_j<\sigma_{ij}.
\]
The associated quadrangle is the union of these two closed order-complex
triangles. It occurs in the ordinary part of the dual manuscript map
exactly when the specified transition ports at $i,j$ are adjacent in the
full boundary cycle of its disk. This construction specifies the ordinary
quadrangles and their attaching maps. The exceptional part consists of
the entire subdivided M\"obius bands. Their rectangular cells are indexed
by the original flags $(v,e)$, and their four vertices are
$U_e^+,S_v^+,S_v^-,U_e^-$ in that cyclic order.

Thus there are specified comparison maps
\[
 \mathcal X\cong\mathcal X^\vee
 \xrightarrow{\ \iota\ }B(M)
 \xrightarrow{\ \beta\ }\overline T(M).
\]
Here $\iota$ is a cellular subcomplex inclusion and induces a surjection on
fundamental groups; $\beta$ is the continuous blow-down collapsing each inserted M\"obius core circle
to its parallel point. The inclusion $\iota$ is a homotopy equivalence
exactly when no ordinary position passes through three old vertices.
In that case it is an identification of the indicated CW spaces.
\end{resultenv}
\begin{proof}
The independent-resolution argument of E2 proves that all $2^\nu$ stable
generizations exist. Cote 154, scan~45, explicitly indexes faces by all
relative positions except the supercritical ones; there is no stable-state
selection in this definition. An explicit position with $\nu=3$ therefore
has eight stable faces, not six. The six-line counterexample and its finite
certificate in Section~\ref{sec:counts} instantiate this discrepancy.

We first identify the subdivision precisely. Give the cube of an ordinary
localization $\sigma$ its coordinates in $[-1,1]^{\nu(\sigma)}$, one for
each zero cocircuit pair. A coordinate remains free on a face if its
zero is unresolved, and equals $+1$ or $-1$ if that zero has been
resolved. Subdivision at coordinate zero therefore gives vertices
indexed by the faces of this original cube. A small $k$-cube is indexed
by an interval of length $k$ in its face order: its least extension has
$k$ specified zeros that may be resolved, and its greatest extension
resolves all of them with the indicated signs. These coordinate
descriptions restrict identically on every common original face, so
the subdivisions glue over the whole ordinary cubical subcomplex.
In particular the edges are exactly cover relations, and the small
two-cubes are exactly incidence diamonds. Triangulating each diamond
along its least-to-greatest diagonal gives its two order simplices.

For the assertion about selection, the two displayed chains are distinct
simplices of the order complex. Their common edge is
$[\sigma,\sigma_{ij}]$, their other edges give the four sides of the
diamond, and their union is an embedded closed disk. In the disk-gluing
construction, a corner bordered by the two indicated transition intervals
is traversed by performing these two independent local resolutions in the
two possible orders. Their supports are disjoint, so the final extension
is the same $\sigma_{ij}$. Conversely an ordinary corner of this type
contains consecutive transition ports, by the definition of a corner of
the disk. This gives the criterion and the attaching map.

No interior simplex is counted twice: a triangle in such a diamond records
its unique least extension, the intermediate resolution, and the final
two-coordinate resolution. The other intermediate vertex is then uniquely
determined. Distinct diamonds therefore meet only in their common faces.
Their union, with their faces, is a simplicial subcomplex after adding the
indicated diagonals. The sign involution carries each diamond to its
opposite diamond, and its quotient remains embedded by the sign argument
in E2.

We now include the exceptional cells, so the description gives the whole
dual surface. Fix an old line $D_i$ and its corresponding parallel
polygon in $\overline T(M)$. Lemma~4.1 of Sturmfels--Ziegler identifies
its boundary with the rank-two circle $\mathcal E(M/i)$, with its order
reversed. After the subdivision just made, denote this boundary circle
by $C_i$. Its $4n_i$ vertices alternate between $S_v^\pm$ and $U_e^\pm$,
and its edges are the ordinary transitions at the flags $v\in\partial e$.
The rank-two sign involution $a_i:C_i\to C_i$ exchanges the two superscripts
and preserves this subdivision. It is a fixed-point-free half-turn.
The interval bundle
\[
 H_i=(C_i\times[-1,1])/\bigl((z,t)\sim(a_i z,-t)\bigr)
\]
is a M\"obius band. Its boundary is identified with $C_i$ by
$z\mapsto[z,1]$. The fiber over the pair $\{z,a_i z\}$ joins these two
boundary vertices. Taking one such fiber edge for every antipodal
vertex pair gives $2n_i$ edges. The complementary closed rectangles,
one per antipodal edge pair of $C_i$, give $2n_i$ two-cells. Each is
the rectangle with corners
$U_e^+,S_v^+,S_v^-,U_e^-$ for one flag $(v,e)$.

The corridor isotopy and endpoint calculation in T3 identify these
fiber edges with the exceptional sweep transitions and these rectangles
with their four-corner circuits. In particular this is equality of
the attaching maps on every boundary edge, not merely a numerical
comparison. Every $C_i$ embeds in the ordinary complex by E2; the band
interiors are disjoint, and each band rectangle is a literal embedded
rectangle in $H_i$. The nonconcurrency hypothesis gives $n_i\geq2$, so
the four vertices of each rectangle are distinct. Consequently the
ordinary selected diamonds together with all the $H_i$ form an embedded
CW subcomplex of the specified $B(M)$. T2 and T3 supply the complete
ordinary and exceptional corner circuits, respectively, so the cell
correspondence is bijective and respects every attaching map. It gives
the asserted cellular homeomorphism from $\mathcal X^\vee$ to that
subcomplex. This also proves that no exceptional cells have been omitted.

The inclusion contains the entire one-skeleton: every ordinary cover
relation occurs as a transition, and all band fiber edges occur as
exceptional transitions. Building $B(M)$ from this subcomplex requires
only the remaining ordinary cells of dimension at least two. Their
attachment induces a quotient on fundamental groups, proving that
$\iota_*$ is surjective. More explicitly, its kernel is the normal
subgroup generated by the boundary loops of the unselected incidence
diamonds: these are precisely the additional two-cells, and attaching
cells of dimension at least three does not change the fundamental group.
Paths from a fixed base point to the attaching loops change their
representatives only by conjugation, so this normal-subgroup description
has no path ambiguity. The additional $k$-cells for $k\geq3$ are all the
ordinary small cubes indexed by length-$k$ face intervals above. In total
the ordinary subdivision has
\[
 N_k=2^k\sum_{j\geq k}\binom jk c_j
     =\frac{2^k}{k!}F^{(k)}(1)
\]
$k$-cells: inside an original open $j$-cube choose its $k$ varying
coordinates and their $2^k$ half-intervals, keeping the other coordinates
zero. Therefore $B(M)$ is obtained from the embedded surface by attaching
$2F''(1)-F'(1)+b$ unselected squares and $N_k$ cells of each dimension
$k\geq3$, where $b=\sum_i n_i$. Here we used T3: the surface has
$F'(1)-b$ ordinary squares and $2b$ band squares. This is a finite
relative CW description of the specified comparison, including every
attaching map through its cube-face interval.

For the blow-down map choose an equivariant
homeomorphism $h_i:C_i\to S^1$, so that
$h_i(a_i z)=-h_i(z)$. In the characteristic disk of the parallel cell
define
\[
 \beta_i([z,t])=t\,h_i(z).
\]
This is well-defined, is the specified boundary identification at
$t=1$, and collapses precisely the core circle at $t=0$. Extend it by
the identity on the ordinary complex. The resulting map $\beta$ has
exactly the stated target and agrees with the ordinary comparison on
every shared edge. Equivalently, $\overline T(M)$ is the quotient of
$B(M)$ by its $n$ disjoint core circles. In particular $\beta_*$ is
surjective and its kernel is the normal subgroup generated by these
core loops, by van Kampen.

Suppose finally that no ordinary position has three zero vertices.
Every ordinary two-cube has exactly four generization ports, no
ordinary specialization port, and no exceptional port. The last claim
follows from T3, because a position near an old line has at most one
zero vertex. T1 makes the four ports alternate between the two zero
vertices, so all four of the subdivided square's diamonds are selected.
There are no higher-dimensional ordinary cubes. Thus the subcomplex
is the whole of $B(M)$, proving the identification in this case. If a
position has at least three zeros, E5 gives a strict Euler-characteristic
difference, so $\iota$ cannot be a homotopy equivalence.

There are $2\nu$ generization ports at a position with $\nu$ zeros. When
there are no intervening specialization ports, their cyclic link has
$2\nu$ adjacent pairs. In general the relevant cycle includes the
specialization ports too; the rule is adjacency in that full cycle. Neither
count is a count of the stable vertices of the $\nu$-cube, which number
$2^\nu$. Furthermore the quadrangles live in its face-incidence subdivision:
their vertices are cube-face centers of dimensions $\nu,\nu-1,\nu-1,\nu-2$.
They are not the two-dimensional faces of the original unsplit cube.
The unmodified extension quotient also does not contain the exceptional
bands by this comparison: its blow-down collapses each entire core
circle. Thus the full cellular embedding belongs to the specified real
blow-up $B(M)$, while its ordinary selected diamonds lie in
$\overline T(M)$ itself.
\end{proof}

\paragraph{Literature scope.}
Sturmfels--Ziegler Theorem~2.3, p.~28, identifies the extension poset of a
\emph{realizable} rank-$r$ oriented matroid with a Baues poset for an
$n$-cube projected to an $(n-r)$-dimensional zonotope. That theorem is an
isomorphism of posets; it neither defines the above cyclic port rule nor
asserts that the manuscript surface is a coherent subcomplex. The rank-four
disconnected examples of Mn\"ev and Richter-Gebert (1993, pp.~271--273)
\cite{MnevRichterGebert1993} do not limit E1--E4, whose rank is exactly three.
Billera, Kapranov, and Sturmfels prove the sphere theorem for cellular
strings of a convex polytope ordered by a generic linear functional
(Theorem~1.2, p.~550), and their coherent-string sphere is the boundary
of the monotone path polytope (\S2, pp.~552--553)
\cite{BilleraKapranovSturmfels1994}. Their higher-dimensional projection
question is explicitly posed as the generalized Baues problem on p.~554.
None of these statements identifies the manuscript's cyclic port rule
with coherence. In particular, the local $\nu$-cube of E2 and the global
$n$-cube in Sturmfels--Ziegler's Baues projection have different roles;
identifying them would not be a type-correct comparison.
We searched the exact titles of the cited papers, ``extension space parallel
extensions punctured'', and ``oriented matroids extension space M\"obius'';
we found no source stating E3, the marked form E4, or the precise cyclic
selector and relative-cell formulas of E6 as used here. Our checked
coverage comprises the complete Sturmfels--Ziegler and cellular-strings
papers, the introduction and first section of Mn\"ev--Richter-Gebert,
and the allowable-sequence definitions of Goodman--Pollack
\cite{GoodmanPollack1980}. The general wiring statements were also
checked in the coauthor chapter, Theorems~5.1.4, 5.2.6, and 5.5.10
\cite{FelsnerGoodman2017}. This is a
report of the search, not a claim of exhaustive coverage or priority;
we supply the full deductions above.
